\uuid{1JNj}
\titre{Matrice/Rang}
\auteur{hanani ; mzari}
\organisation{exo7}
\licence{CC BY-NC-SA 4.0}
\datecreate{2020-01-20}
\niveau{L1}
\module{Algèbre}
\chapitre{Matrice}
\sousChapitre{Propriétés élémentaires, généralités}
\difficulte{3}
\theme{Calcul matriciel, Matrice, Rang}
\exo7qcmid{lille-2.matrices.013}

\contenu{
\begin{question}
\qclassification{108.01}
\qtags{motcle=Matrice/Rang}
Soit $A$ une matrice de rang $r$. Quelles sont les assertions vraies ?
\begin{answers}
\good{$A$ admet $r$ vecteurs colonnes linéairement indépendants}.
\good{$A$ admet $r$ vecteurs lignes linéairement indépendants}.
\bad{Toute famille contenant $r$ vecteurs colonnes de $A$ est libre}.
\bad{Toute famille contenant $r$ vecteurs lignes de $A$ est libre}.
\end{answers}
\begin{explanations} Le rang d'une matrice est le nombre maximum de vecteurs colonnes ou  lignes qui sont linéairement indépendants.
\end{explanations}
\end{question}
}
